\chapter{Rubrik Penilaian Lengkap Berbasis OBE}
\label{app:rubrik}

Lampiran ini memuat seluruh rubrik yang dipakai pada mata kuliah ini.
Setiap rubrik ditautkan secara eksplisit ke CPMK sehingga penilaian
dapat dipertanggungjawabkan.

\section{Rekapitulasi Bobot Penilaian}

\begin{table}[H]
    \centering
    \caption{Komponen penilaian dan bobotnya.}
    \label{tab:app-bobot}
    \small
    \begin{tabular}{@{}p{3.4cm}p{4.2cm}p{3.2cm}c@{}}
        \toprule
        \textbf{Komponen} & \textbf{Teknik asesmen} &
        \textbf{CPMK diukur} & \textbf{Bobot} \\
        \midrule
        Tugas Individu & Analisis dan pemrograman &
        \CPMK{01}, \CPMK{02} & 15\% \\
        \addlinespace
        Praktikum & Latihan laboratorium &
        \CPMK{03}, \CPMK{04} & 20\% \\
        \addlinespace
        Kuis & Teori dan praktik &
        \CPMK{01}, \CPMK{02} & 10\% \\
        \addlinespace
        UTS (minggu 8) & Ujian tertulis dan \eng{coding} &
        \CPMK{01}, \CPMK{02} & 20\% \\
        \addlinespace
        Proyek Akhir & Aplikasi akhir &
        \CPMK{03}, \CPMK{04} & 15\% \\
        \addlinespace
        UAS (minggu 16) & Ujian tertulis dan \eng{coding} &
        \CPMK{01}--\CPMK{04} & 20\% \\
        \midrule
        \textbf{Total} & & & \textbf{100\%} \\
        \bottomrule
    \end{tabular}
\end{table}

\subsection{Kriteria Konversi Nilai}

\begin{table}[H]
    \centering
    \caption{Kriteria dan deskripsi tingkat penguasaan.}
    \label{tab:app-konversi}
    \small
    \begin{tabular}{@{}ccp{8.4cm}@{}}
        \toprule
        \textbf{Nilai} & \textbf{Rentang} & \textbf{Deskripsi} \\
        \midrule
        A & 85--100 & Menguasai semua CPMK dengan sangat baik;
        mampu mengembangkan program assembly kompleks secara mandiri \\
        \addlinespace
        B & 70--84 & Menguasai sebagian besar CPMK dengan baik;
        mampu membuat program assembly moderat \\
        \addlinespace
        C & 60--69 & Menguasai CPMK dasar dengan cukup; mampu
        membuat program assembly sederhana \\
        \addlinespace
        D & 50--59 & Menguasai sebagian kecil CPMK; program assembly
        terbatas \\
        \addlinespace
        E & $<$50 & Belum menguasai CPMK yang ditetapkan \\
        \bottomrule
    \end{tabular}
\end{table}

\section{Rubrik Tugas Individu}

\begin{table}[H]
    \centering
    \caption{Rubrik penilaian tugas individu (Tugas 1--7).}
    \label{tab:app-rubrik-tugas}
    \footnotesize
    \begin{tabular}{@{}p{3.0cm}p{3.4cm}p{3.4cm}p{3.4cm}@{}}
        \toprule
        \textbf{Kriteria} & \textbf{4 -- Sangat baik} &
        \textbf{3 -- Baik} & \textbf{2 -- Cukup} \\
        \midrule
        Ketepatan jawaban konseptual (\CPMK{01}, \CPMK{02}) &
        Seluruh jawaban benar dengan penjelasan sebab-akibat &
        Sebagian besar benar, penjelasan singkat &
        Jawaban benar sebagian, tanpa penjelasan \\
        \addlinespace
        Ketepatan perhitungan (\CPMK{01}) & Semua perhitungan benar
        dan langkah ditulis & Semua benar, langkah sebagian &
        Ada kesalahan perhitungan \\
        \addlinespace
        Kualitas program (\CPMK{03}) & Program berjalan benar,
        terstruktur, dan berkomentar & Program berjalan benar &
        Program berjalan dengan koreksi kecil \\
        \addlinespace
        Kelengkapan dokumentasi & Laporan lengkap dengan tangkapan
        layar bukti & Laporan lengkap tanpa bukti visual &
        Laporan kurang lengkap \\
        \bottomrule
    \end{tabular}
\end{table}

\section{Rubrik Praktikum}

\begin{table}[H]
    \centering
    \caption{Rubrik penilaian praktikum.}
    \label{tab:app-rubrik-praktikum}
    \footnotesize
    \begin{tabular}{@{}p{3.0cm}p{3.4cm}p{3.4cm}p{3.4cm}@{}}
        \toprule
        \textbf{Kriteria} & \textbf{4 -- Sangat baik} &
        \textbf{3 -- Baik} & \textbf{2 -- Cukup} \\
        \midrule
        Kebenaran hasil praktikum (\CPMK{03}) & Semua langkah
        berhasil, keluaran benar & Langkah utama berhasil &
        Berhasil dengan bantuan \\
        \addlinespace
        Keterampilan alat (\CPMK{04}) & Menggunakan \eng{debugger}
        dan alat bantu secara mandiri dan efisien &
        Menggunakan dengan lancar & Menggunakan dengan panduan \\
        \addlinespace
        Kualitas jurnal & Jurnal lengkap, analitis, dan jujur &
        Jurnal lengkap & Jurnal kurang lengkap \\
        \addlinespace
        Kerja sama (bila berpasangan) & Kontribusi setara dan aktif
        membantu & Kontribusi baik & Kontribusi terbatas \\
        \bottomrule
    \end{tabular}
\end{table}

\section{Rubrik Kuis}

Kuis dinilai secara objektif. Setiap butir memiliki bobot yang sama
di dalam kuis yang bersangkutan, kecuali dinyatakan lain pada
kisi-kisi bab. Nilai kuis adalah persentase jawaban benar. Tidak ada
penilaian sebagian untuk soal pilihan ganda.

\section{Rubrik Proyek Akhir}

Rubrik proyek akhir disajikan pada Tabel~\ref{tab:rubrik-proyek} di
Bab~\ref{chap:evaluasi-akhir}, dengan lima kriteria berbobot sama:
kebenaran fungsi, kualitas kode, penggunaan instruksi, pengujian dan
\eng{debugging}, serta laporan dan presentasi.

\section{Rubrik Ujian Tertulis dan \eng{Coding}}

\begin{table}[H]
    \centering
    \caption{Rubrik penilaian soal \eng{coding} pada UTS dan UAS.}
    \label{tab:app-rubrik-coding}
    \footnotesize
    \begin{tabular}{@{}p{3.0cm}p{3.4cm}p{3.4cm}p{3.4cm}@{}}
        \toprule
        \textbf{Kriteria} & \textbf{4 -- Sangat baik} &
        \textbf{3 -- Baik} & \textbf{2 -- Cukup} \\
        \midrule
        Kebenaran logika & Algoritma benar dan lengkap &
        Algoritma benar, ada detail terlewat & Algoritma benar
        sebagian \\
        \addlinespace
        Ketepatan sintaks & Tidak ada kesalahan sintaks &
        Satu-dua kesalahan kecil & Beberapa kesalahan sintaks \\
        \addlinespace
        Penggunaan instruksi (\CPMK{02}) & Instruksi dan mode
        pengalamatan dipilih tepat & Sebagian besar tepat &
        Banyak pilihan tidak tepat \\
        \addlinespace
        Penanganan kasus batas (\CPMK{04}) & Semua kasus batas
        ditangani & Sebagian kasus batas ditangani &
        Kasus batas diabaikan \\
        \bottomrule
    \end{tabular}
\end{table}

\section{Pemetaan Rubrik ke CPMK}

\begin{table}[H]
    \centering
    \caption{Pemetaan setiap rubrik ke CPMK yang diukurnya.}
    \label{tab:app-pemetaan-rubrik}
    \small
    \begin{tabular}{@{}p{4.0cm}p{2.2cm}p{2.2cm}p{2.2cm}p{2.2cm}@{}}
        \toprule
        \textbf{Rubrik} & \textbf{\CPMK{01}} & \textbf{\CPMK{02}} &
        \textbf{\CPMK{03}} & \textbf{\CPMK{04}} \\
        \midrule
        Tugas Individu & \ding{51} & \ding{51} & \ding{51} & \\
        \addlinespace
        Praktikum & & & \ding{51} & \ding{51} \\
        \addlinespace
        Kuis & \ding{51} & \ding{51} & & \\
        \addlinespace
        UTS & \ding{51} & \ding{51} & & \\
        \addlinespace
        Proyek Akhir & & \ding{51} & \ding{51} & \ding{51} \\
        \addlinespace
        UAS & \ding{51} & \ding{51} & \ding{51} & \ding{51} \\
        \bottomrule
    \end{tabular}
\end{table}
